\chapter{Testing Philosophy}
\label{ch:testing-philosophy}

Almost every specific number and bug this book has cited across both volumes came from one
of a small number of testing mechanisms. This chapter gathers them into one place, because
seen individually they are backend-specific facts, but seen together they are a coherent
testing philosophy worth naming explicitly.

\section{The base layer: GoogleTest, at real scale}

CNA's own unit and integration suite (\texttt{CnaTests}) runs to roughly 4{,}370 tests at the
scale referenced throughout this book, plus a separate GPU pixel-test suite of roughly 490
tests across the four backends mature enough to support it. \texttt{sharp-runtime}, the
foundation layer covered in Chapter~\ref{ch:sharp-runtime-overview}, runs a suite of its own
numbering in the thousands. None of this is exotic --- GoogleTest, run through CTest --- but
the scale matters for what it makes possible: a change that regresses one backend's behavior
without touching another's has somewhere to actually surface.

\section{Pixel tests: the level where most of this book's bugs were actually found}

An ordinary unit test can confirm a method returns without crashing, or that a getter returns
what a setter most recently stored --- neither can catch that
\texttt{transformMatrix} was silently ignored by \cnaclass{SpriteBatch::Begin()}
(Chapter~\ref{ch:sdl-renderer}), or that a Vulkan blend-state implementation applied one
hardcoded equation regardless of what was requested (Chapter~\ref{ch:vulkan-backend}). Both
of those --- and the substantial majority of the specific, named bugs in this book's Part~IV
and Part~III chapters --- were only found by a pixel test: rendering a real frame and reading
real pixel values back, then asserting on the actual color, not merely on the absence of a
crash. This is the level at which CNA's own house style Chapter~\ref{ch:design-philosophy}
described bites hardest: a stub that silently drops a parameter compiles cleanly, passes a
crash-only test, and only fails once something reads the pixels it was supposed to affect.

\section{Differential testing: a real reference, not a specification}

Chapter~\ref{ch:what-is-cna} named this as one of CNA's three most distinctive verification
mechanisms, and it is worth restating precisely here: \texttt{tools/fna-reference/} compares
CNA's behavior against a genuinely running FNA installation, not against a written
description of what FNA is supposed to do. This distinction mattered concretely in
Chapter~\ref{ch:media-system}'s \cnaclass{Song} finding --- eight audit rounds diffed against
FNA's own source and missed a real gap, because FNA's source was itself missing something
relative to the true XNA~4.0 reference. A specification, however carefully read, only tells
you what its author believed to be true; a running reference implementation tells you what
actually happens.

\section{The oracle: pixel-for-pixel against real XNA itself}

\texttt{D3D9}'s own verification, covered in full in Chapter~\ref{ch:direct3d-backends}, goes
one level further still: not a comparison against FNA, but against a genuine, running XNA~4.0
runtime, diffed at zero tolerance across a corpus that had grown to 36 scenes with zero
divergence at the time of writing. This is the strictest bar in the entire project, and it is
worth remembering that it caught two real bugs neither differential-against-FNA testing nor
ordinary pixel tests had caught: the sort-mode Z-clipping bug and the
\cnaclass{GraphicsProfile}-never-reaches-the-device bug, both covered in
Chapter~\ref{ch:direct3d-backends}.

\section{Compile-time purity: NOXNA as a build gate}

Not every verification mechanism in this project runs at test time. The \texttt{NOXNA} build
option (Chapter~\ref{ch:design-philosophy}) turns every untagged non-XNA declaration inside
the XNA namespace tree into a \texttt{[[deprecated]]} warning under \texttt{-Werror} ---
catching a specific class of drift (a project-specific extension silently leaking into the
public, supposedly XNA-faithful surface) at compile time, before a single test ever runs.

\section{Sanitizers, and where they are used}

Memory-safety and undefined-behavior sanitizers are used deliberately, not universally.
\texttt{free-direct} (Chapter~\ref{ch:freedirect-deep-dive}) has a dedicated, opt-in sanitizer
build configuration verified clean across its own full test suite, including in combination
with its optional ENet networking transport. \cnans{Microsoft::Devices::Sensors}
(Chapter~\ref{ch:devices-sensors}) had a genuine data race found by exactly this kind of
tooling --- an unguarded read of a disposal flag from a background sensor-event thread --- and
a later signed-integer-overflow bug in its own update-throttling logic was confirmed
specifically under UndefinedBehaviorSanitizer, not by inspection. Sanitizer coverage in this
project is targeted at the specific subsystems where genuine concurrency or native-interop
risk exists, rather than blanket-applied everywhere.

\section{What CI actually covers, stated honestly}

Chapter~\ref{ch:building-cna} already gave this fact once, and it is worth
repeating here in its proper context: automated CI for this project is Linux-only and
partial, running only the Input and Devices/Sensors suites at the time of writing --- not the
full unit suite, and not the GPU pixel-test suite across any backend. Every Graphics-layer
claim in this entire book rests on suites run locally and by hand, not on a green CI badge.
There is no Windows, macOS, or Android CI at all; the Windows-targeting backends in
Chapter~\ref{ch:direct3d-backends} are verified through the Wine-based dev loop covered in
this Part's cross-platform chapters, and Android is verified against a real emulator by hand,
not automatically on every commit. This is not a criticism the project hides --- it is stated
plainly in its own README --- and it is the reason this book has consistently distinguished
``verified locally, by a specific mechanism, on a specific date'' from ``covered by CI,''
rather than treating the two as interchangeable.

\section{The throughline}

Every mechanism in this chapter answers the same underlying question from a different angle:
\emph{how do we know this is actually true, rather than merely believed to be true?} A unit
test answers it for logic in isolation. A pixel test answers it for what actually reaches the
screen. Differential testing against FNA answers it against a real, if imperfect, reference.
The D3D9 oracle answers it against the most authoritative reference available at all. The
\texttt{NOXNA} build gate answers a narrower, structural version of the question before any
of the others even run. None of these mechanisms alone would have caught every bug this book
has named across both volumes --- it is the combination, applied unevenly and honestly
according to what each subsystem actually needed, that this chapter's title refers to as a
philosophy rather than a single technique.
